\chapter{Hadronenmultipletts}

\section{Ziel dieses Kapitels}

Dieses Kapitel behandelt die Ordnung von Hadronen in Multipletts.
Hadronenmultipletts zeigen, dass Hadronen nicht zufaellig auftreten, sondern nach
Quantenzahlen wie Isospin, Hyperladung, Strangeness und Spin systematisch geordnet
werden koennen.

\begin{notebox}
Dieses Kapitel behandelt die Stichwoerter:
\begin{itemize}
    \item Hadronenmultipletts,
    \item Isospin,
    \item dritte Isospinkomponente $I_3$,
    \item Hyperladung $Y$,
    \item Strangeness $S$,
    \item Baryonenoktett,
    \item Baryonendekuplett,
    \item Mesonenoktett,
    \item Mesonennonett,
    \item Pionen,
    \item Kaonen,
    \item Nukleonen,
    \item Delta-Teilchen,
    \item Omega-Baryon,
    \item Quarkinhalt von Hadronen.
\end{itemize}
\end{notebox}

\section{Warum Multipletts?}

Hadronen besitzen viele verschiedene Ladungen, Massen und Zerfallseigenschaften.
Trotzdem erkennt man Muster.

\begin{conceptbox}
Hadronen mit aehnlichen Eigenschaften lassen sich in Gruppen anordnen.
Diese Gruppen nennt man Multipletts.

Die Ordnung erfolgt typischerweise nach:
\begin{itemize}
    \item Isospin $I$,
    \item dritter Isospinkomponente $I_3$,
    \item Hyperladung $Y$,
    \item Strangeness $S$,
    \item Spin $J$,
    \item Quarkinhalt.
\end{itemize}
\end{conceptbox}

\begin{definitionbox}
Ein \emph{Hadronenmultiplett} ist eine Gruppe von Hadronen, die durch gemeinsame
Quantenzahlen und Symmetrien zusammengehoeren.
\end{definitionbox}

\begin{examtipbox}
Multipletts sind eine Ordnungshilfe.

Man soll erkennen koennen:
\begin{itemize}
    \item Welche Teilchen gehoeren zusammen?
    \item Welche Quantenzahlen unterscheiden sie?
    \item Welchen Quarkinhalt haben sie?
    \item Wie berechnet man ihre Ladung?
\end{itemize}
\end{examtipbox}

\section{Achsen eines Multiplettdiagramms}

Hadronenmultipletts werden oft in einem Diagramm dargestellt.
Die horizontale Achse ist meist die dritte Isospinkomponente $I_3$.
Die vertikale Achse ist meist die Hyperladung $Y$.

\begin{definitionbox}
Die \emph{dritte Isospinkomponente} $I_3$ unterscheidet die Ladungszustaende innerhalb
eines Isospinmultipletts.
\end{definitionbox}

\begin{definitionbox}
Die \emph{Hyperladung} ist im einfachen strange-Hadron-Kontext:
\[
    Y = B + S.
\]
Dabei ist $B$ die Baryonenzahl und $S$ die Strangeness.
\end{definitionbox}

\begin{formulabox}
Die elektrische Ladung folgt aus der Gell-Mann--Nishijima-Formel:
\[
    Q = I_3 + \frac{Y}{2}.
\]
\end{formulabox}

\begin{conceptbox}
In einem Multiplettdiagramm kann man die elektrische Ladung aus der Position
ablesen, wenn $I_3$ und $Y$ bekannt sind.

Teilchen mit gleicher Hyperladung liegen auf derselben horizontalen Linie.
Teilchen mit unterschiedlichen Werten von $I_3$ liegen nebeneinander.
\end{conceptbox}

\section{Quarkinhalt und Quantenzahlen}

Die wichtigsten leichten Quarks fuer diese Multipletts sind:
\[
    u,\qquad d,\qquad s.
\]

\begin{notebox}
Quarkladungen:
\[
    Q(u)=+\frac{2}{3},
    \qquad
    Q(d)=-\frac{1}{3},
    \qquad
    Q(s)=-\frac{1}{3}.
\]

Strangeness:
\[
    S(u)=0,
    \qquad
    S(d)=0,
    \qquad
    S(s)=-1.
\]
\end{notebox}

\begin{notebox}
Antiquarks haben entgegengesetzte additive Quantenzahlen:
\[
    Q(\overline{u})=-\frac{2}{3},
    \qquad
    Q(\overline{d})=+\frac{1}{3},
    \qquad
    Q(\overline{s})=+\frac{1}{3}.
\]
\[
    S(\overline{s})=+1.
\]
\end{notebox}

\begin{examtipbox}
Bei jedem Hadron zuerst den Quarkinhalt aufschreiben.
Dann addiert man:
\begin{itemize}
    \item elektrische Ladung $Q$,
    \item Baryonenzahl $B$,
    \item Strangeness $S$,
    \item Hyperladung $Y=B+S$.
\end{itemize}
\end{examtipbox}

\section{Baryonen und Mesonen}

\begin{definitionbox}
\emph{Baryonen} bestehen im einfachen Quarkmodell aus drei Quarks:
\[
    qqq.
\]
Sie haben Baryonenzahl
\[
    B=1.
\]
\end{definitionbox}

\begin{definitionbox}
\emph{Mesonen} bestehen im einfachen Quarkmodell aus einem Quark und einem Antiquark:
\[
    q\overline{q}.
\]
Sie haben Baryonenzahl
\[
    B=0.
\]
\end{definitionbox}

\begin{conceptbox}
Da Baryonen und Mesonen unterschiedliche Baryonenzahl besitzen, werden sie in
getrennten Multipletts angeordnet.
\end{conceptbox}

\section{Nukleonen als Isospin-Dublett}

Proton und Neutron bilden ein Isospin-Dublett.

\begin{definitionbox}
Ein \emph{Isospin-Dublett} besitzt
\[
    I=\frac{1}{2}
\]
und damit zwei Werte von $I_3$:
\[
    I_3=+\frac{1}{2},
    \qquad
    I_3=-\frac{1}{2}.
\]
\end{definitionbox}

\begin{conceptbox}
Fuer die Nukleonen gilt:
\[
    p=uud,
    \qquad
    n=udd.
\]

Beide haben:
\[
    B=1,
    \qquad
    S=0,
    \qquad
    Y=B+S=1.
\]
\end{conceptbox}

\begin{center}
\begin{tabular}{|c|c|c|c|c|}
\hline
Teilchen & Quarkinhalt & $I_3$ & $Y$ & $Q$ \\
\hline
$p$ & $uud$ & $+\frac{1}{2}$ & $1$ & $+1$ \\
\hline
$n$ & $udd$ & $-\frac{1}{2}$ & $1$ & $0$ \\
\hline
\end{tabular}
\end{center}

\begin{examplebox}
Fuer das Proton gilt:
\[
    I_3=+\frac{1}{2},
    \qquad
    Y=1.
\]
Mit
\[
    Q=I_3+\frac{Y}{2}
\]
folgt:
\[
    Q=+\frac{1}{2}+\frac{1}{2}=+1.
\]
\end{examplebox}

\begin{examplebox}
Fuer das Neutron gilt:
\[
    I_3=-\frac{1}{2},
    \qquad
    Y=1.
\]
Also:
\[
    Q=-\frac{1}{2}+\frac{1}{2}=0.
\]
\end{examplebox}

\section{Pionen als Isospin-Triplett}

Die Pionen bilden ein Isospin-Triplett.

\begin{definitionbox}
Ein \emph{Isospin-Triplett} besitzt
\[
    I=1
\]
und damit drei Werte von $I_3$:
\[
    I_3=+1,\quad 0,\quad -1.
\]
\end{definitionbox}

\begin{conceptbox}
Die Pionen sind Mesonen.
Sie haben:
\[
    B=0,
    \qquad
    S=0,
    \qquad
    Y=0.
\]
\end{conceptbox}

\begin{center}
\begin{tabular}{|c|c|c|c|c|}
\hline
Teilchen & Quarkinhalt & $I_3$ & $Y$ & $Q$ \\
\hline
$\pi^+$ & $u\overline{d}$ & $+1$ & $0$ & $+1$ \\
\hline
$\pi^0$ & Mischung aus $u\overline{u}$ und $d\overline{d}$ & $0$ & $0$ & $0$ \\
\hline
$\pi^-$ & $d\overline{u}$ & $-1$ & $0$ & $-1$ \\
\hline
\end{tabular}
\end{center}

\begin{examtipbox}
Bei Pionen gilt:
\[
    Y=0.
\]
Daher ist nach
\[
    Q=I_3+\frac{Y}{2}
\]
einfach:
\[
    Q=I_3.
\]
\end{examtipbox}

\section{Kaonen und Antikaonen}

Kaonen enthalten ein strange-Quark oder ein strange-Antiquark.

\begin{conceptbox}
Kaonen mit einem strange-Antiquark haben positive Strangeness:
\[
    S(\overline{s})=+1.
\]

Antikaonen mit einem strange-Quark haben negative Strangeness:
\[
    S(s)=-1.
\]
\end{conceptbox}

\begin{center}
\begin{tabular}{|c|c|c|c|c|}
\hline
Teilchen & Quarkinhalt & $S$ & $Y$ & $Q$ \\
\hline
$K^+$ & $u\overline{s}$ & $+1$ & $+1$ & $+1$ \\
\hline
$K^0$ & $d\overline{s}$ & $+1$ & $+1$ & $0$ \\
\hline
$\overline{K}^0$ & $s\overline{d}$ & $-1$ & $-1$ & $0$ \\
\hline
$K^-$ & $s\overline{u}$ & $-1$ & $-1$ & $-1$ \\
\hline
\end{tabular}
\end{center}

\begin{examtipbox}
Vorzeichen der Strangeness merken:
\[
    S(s)=-1,
    \qquad
    S(\overline{s})=+1.
\]

Deshalb haben $K^+$ und $K^0$ positive Strangeness, waehrend $K^-$ und
$\overline{K}^0$ negative Strangeness haben.
\end{examtipbox}

\section{Mesonenoktett und Mesonennonett}

Mesonen aus den leichten Quarks $u$, $d$ und $s$ koennen in Multipletts angeordnet
werden.

\begin{definitionbox}
Das \emph{Mesonenoktett} ist ein Multiplett aus acht Mesonen.
Zusammen mit einem Singulett entsteht ein \emph{Mesonennonett}.
\end{definitionbox}

\begin{conceptbox}
Fuer Spin-$0$-Mesonen gehoeren typischerweise dazu:
\begin{itemize}
    \item Pionen $\pi^+$, $\pi^0$, $\pi^-$,
    \item Kaonen $K^+$, $K^0$, $\overline{K}^0$, $K^-$,
    \item eta-artige neutrale Mesonen.
\end{itemize}
\end{conceptbox}

\begin{notebox}
In den Vorlesungsfolien wird ein Mesonennonett fuer Spin-$0$-Mesonen gezeigt.
Darin treten Pionen, Kaonen, Antikaonen, Eta- und Eta-Strich-Mesonen auf.
\end{notebox}

\begin{conceptbox}
Fuer Spin-$1$-Mesonen gibt es ein entsprechendes Vektormesonen-Nonett.

Darin treten zum Beispiel rho-artige, omega-artige, phi-artige und angeregte
Kaon-Zustaende auf.
\end{conceptbox}

\begin{examtipbox}
Nonett bedeutet:
\[
    8+1=9.
\]

Das Oktett kommt aus der symmetrischen Ordnung der leichten Quarkflavors.
Das zusaetzliche Singulett mischt sich bei realen neutralen Mesonen mit den
Oktettzustaenden.
\end{examtipbox}

\section{Baryonenoktett}

Das Baryonenoktett enthaelt Spin-$\frac{1}{2}$-Baryonen aus den leichten Quarks
$u$, $d$ und $s$.

\begin{definitionbox}
Das \emph{Baryonenoktett} ist ein Multiplett aus acht Baryonen mit Spin
\[
    J=\frac{1}{2}.
\]
\end{definitionbox}

\begin{conceptbox}
Zum Baryonenoktett gehoeren:
\begin{itemize}
    \item Nukleonen $p$ und $n$,
    \item Sigma-Baryonen $\Sigma^+$, $\Sigma^0$, $\Sigma^-$,
    \item Lambda-Baryon $\Lambda^0$,
    \item Xi-Baryonen $\Xi^0$ und $\Xi^-$.
\end{itemize}
\end{conceptbox}

\begin{center}
\begin{tabular}{|c|c|c|c|c|}
\hline
Teilchen & Quarkinhalt & $S$ & $Y=B+S$ & $Q$ \\
\hline
$p$ & $uud$ & $0$ & $1$ & $+1$ \\
\hline
$n$ & $udd$ & $0$ & $1$ & $0$ \\
\hline
$\Sigma^+$ & $uus$ & $-1$ & $0$ & $+1$ \\
\hline
$\Sigma^0$ & $uds$ & $-1$ & $0$ & $0$ \\
\hline
$\Sigma^-$ & $dds$ & $-1$ & $0$ & $-1$ \\
\hline
$\Lambda^0$ & $uds$ & $-1$ & $0$ & $0$ \\
\hline
$\Xi^0$ & $uss$ & $-2$ & $-1$ & $0$ \\
\hline
$\Xi^-$ & $dss$ & $-2$ & $-1$ & $-1$ \\
\hline
\end{tabular}
\end{center}

\begin{notebox}
$\Sigma^0$ und $\Lambda^0$ haben beide Quarkinhalt $uds$, unterscheiden sich aber
in Isospin und innerer Symmetrie.

Deshalb sind sie verschiedene Teilchen.
\end{notebox}

\begin{examtipbox}
Das Baryonenoktett hat Spin
\[
    J=\frac{1}{2}.
\]

Die wichtigsten Mitglieder sind:
\[
    N,\quad \Lambda,\quad \Sigma,\quad \Xi.
\]
\end{examtipbox}

\section{Baryonendekuplett}

Das Baryonendekuplett enthaelt Spin-$\frac{3}{2}$-Baryonen.

\begin{definitionbox}
Das \emph{Baryonendekuplett} ist ein Multiplett aus zehn Baryonen mit Spin
\[
    J=\frac{3}{2}.
\]
\end{definitionbox}

\begin{conceptbox}
Zum Baryonendekuplett gehoeren:
\begin{itemize}
    \item Delta-Teilchen $\Delta^{++}$, $\Delta^+$, $\Delta^0$, $\Delta^-$,
    \item angeregte Sigma-Baryonen $\Sigma^{*+}$, $\Sigma^{*0}$, $\Sigma^{*-}$,
    \item angeregte Xi-Baryonen $\Xi^{*0}$, $\Xi^{*-}$,
    \item Omega-Baryon $\Omega^-$.
\end{itemize}
\end{conceptbox}

\begin{center}
\begin{tabular}{|c|c|c|c|c|}
\hline
Teilchen & Quarkinhalt & $S$ & $Y=B+S$ & $Q$ \\
\hline
$\Delta^{++}$ & $uuu$ & $0$ & $1$ & $+2$ \\
\hline
$\Delta^{+}$ & $uud$ & $0$ & $1$ & $+1$ \\
\hline
$\Delta^{0}$ & $udd$ & $0$ & $1$ & $0$ \\
\hline
$\Delta^{-}$ & $ddd$ & $0$ & $1$ & $-1$ \\
\hline
$\Sigma^{*+}$ & $uus$ & $-1$ & $0$ & $+1$ \\
\hline
$\Sigma^{*0}$ & $uds$ & $-1$ & $0$ & $0$ \\
\hline
$\Sigma^{*-}$ & $dds$ & $-1$ & $0$ & $-1$ \\
\hline
$\Xi^{*0}$ & $uss$ & $-2$ & $-1$ & $0$ \\
\hline
$\Xi^{*-}$ & $dss$ & $-2$ & $-1$ & $-1$ \\
\hline
$\Omega^{-}$ & $sss$ & $-3$ & $-2$ & $-1$ \\
\hline
\end{tabular}
\end{center}

\begin{examtipbox}
Das Omega-Baryon ist besonders wichtig:
\[
    \Omega^- = sss.
\]
Daraus folgt:
\[
    S=-3,
    \qquad
    Q=-1.
\]
\end{examtipbox}

\section{Delta-Teilchen}

Delta-Teilchen sind angeregte Baryonen mit Spin
\[
    J=\frac{3}{2}.
\]

\begin{conceptbox}
Die Delta-Teilchen bilden ein Isospin-Quartett:
\[
    \Delta^{++},
    \qquad
    \Delta^+,
    \qquad
    \Delta^0,
    \qquad
    \Delta^-.
\]

Sie haben:
\[
    I=\frac{3}{2}.
\]
\end{conceptbox}

\begin{center}
\begin{tabular}{|c|c|c|}
\hline
Teilchen & Quarkinhalt & Ladung \\
\hline
$\Delta^{++}$ & $uuu$ & $+2$ \\
\hline
$\Delta^+$ & $uud$ & $+1$ \\
\hline
$\Delta^0$ & $udd$ & $0$ \\
\hline
$\Delta^-$ & $ddd$ & $-1$ \\
\hline
\end{tabular}
\end{center}

\begin{conceptbox}
$\Delta^+$ und Proton haben denselben Quarkinhalt:
\[
    uud.
\]

Sie unterscheiden sich aber im Spin und im inneren Zustand.
Das Proton gehoert zum Spin-$\frac{1}{2}$-Baryonenoktett.
Das $\Delta^+$ gehoert zum Spin-$\frac{3}{2}$-Baryonendekuplett.
\end{conceptbox}

\begin{examtipbox}
Gleicher Quarkinhalt bedeutet nicht automatisch gleiches Teilchen.
Spin, Isospin und innerer Zustand sind ebenfalls entscheidend.
\end{examtipbox}

\section{Omega-Baryon}

Das Omega-Minus-Baryon ist ein besonders klares Mitglied des Dekupletts.

\begin{definitionbox}
Das Omega-Minus-Baryon hat Quarkinhalt:
\[
    \Omega^- = sss.
\]
\end{definitionbox}

\begin{conceptbox}
Fuer das Omega-Minus gilt:
\[
    Q
    =
    -\frac{1}{3}
    -
    \frac{1}{3}
    -
    \frac{1}{3}
    =
    -1.
\]

Da drei strange-Quarks enthalten sind:
\[
    S=-3.
\]

Da es ein Baryon ist:
\[
    B=1.
\]

Also:
\[
    Y=B+S=1-3=-2.
\]
\end{conceptbox}

\begin{examtipbox}
$\Omega^-$ ist pruefungsrelevant, weil es die unterste Spitze des
Baryonendekupletts bildet:
\[
    S=-3,
    \qquad
    Y=-2,
    \qquad
    Q=-1.
\]
\end{examtipbox}

\section{Ladungsberechnung aus Quarkinhalten}

\begin{examplebox}
Berechne die Ladung von $\Sigma^+$:
\[
    \Sigma^+ = uus.
\]
Dann:
\[
    Q
    =
    \frac{2}{3}
    +
    \frac{2}{3}
    -
    \frac{1}{3}
    =
    +1.
\]
\end{examplebox}

\begin{examplebox}
Berechne die Ladung von $\Xi^-$:
\[
    \Xi^- = dss.
\]
Dann:
\[
    Q
    =
    -\frac{1}{3}
    -
    \frac{1}{3}
    -
    \frac{1}{3}
    =
    -1.
\]
\end{examplebox}

\begin{examplebox}
Berechne die Ladung von $K^0$:
\[
    K^0=d\overline{s}.
\]
Dann:
\[
    Q
    =
    -\frac{1}{3}
    +
    \frac{1}{3}
    =
    0.
\]
\end{examplebox}

\section{Multipletts und starke Wechselwirkung}

\begin{conceptbox}
Hadronenmultipletts sind eng mit der starken Wechselwirkung verbunden.

Die starke Wechselwirkung behandelt die leichten Quarkflavors $u$, $d$ und $s$
naeherungsweise symmetrisch.
Diese Symmetrie ist nicht exakt, weil das strange-Quark schwerer ist als $u$ und $d$.
\end{conceptbox}

\begin{notebox}
Wenn die Symmetrie exakt waere, haetten Teilchen innerhalb eines Multipletts gleiche
Massen.

In Wirklichkeit sind die Massen verschieden.
Die Massenunterschiede entstehen unter anderem durch unterschiedliche Quarkmassen
und elektromagnetische Effekte.
\end{notebox}

\begin{examtipbox}
Multipletts zeigen:
Teilchen mit unterschiedlichen Ladungen und Strangeness koennen zur gleichen
symmetrischen Familie gehoeren.
\end{examtipbox}

\section{Massenmuster in Multipletts}

\begin{conceptbox}
Innerhalb eines Multipletts steigen die Massen typischerweise, wenn mehr strange-Quarks
enthalten sind.

Grund:
Das strange-Quark ist schwerer als up- und down-Quarks.
\end{conceptbox}

\begin{notebox}
Typische qualitative Reihenfolge im Baryonenoktett:
\[
    N
    \quad < \quad
    \Lambda,\Sigma
    \quad < \quad
    \Xi.
\]

Typische qualitative Reihenfolge im Baryonendekuplett:
\[
    \Delta
    \quad < \quad
    \Sigma^*
    \quad < \quad
    \Xi^*
    \quad < \quad
    \Omega^-.
\]
\end{notebox}

\begin{examtipbox}
Mehr strange-Quarks bedeuten meistens groessere Masse.
Das ist ein wichtiger qualitativer Trend in den Multipletts.
\end{examtipbox}

\section{Darstellung im $I_3$-$Y$-Diagramm}

\begin{conceptbox}
In einem $I_3$-$Y$-Diagramm gilt:
\begin{itemize}
    \item horizontale Richtung: $I_3$,
    \item vertikale Richtung: $Y$,
    \item elektrische Ladung: $Q=I_3+Y/2$.
\end{itemize}
\end{conceptbox}

\begin{examplebox}
Betrachte eine horizontale Linie mit
\[
    Y=1.
\]
Dann gilt:
\[
    Q=I_3+\frac{1}{2}.
\]

Fuer
\[
    I_3=+\frac{1}{2}
\]
folgt:
\[
    Q=+1.
\]

Fuer
\[
    I_3=-\frac{1}{2}
\]
folgt:
\[
    Q=0.
\]

Das passt zu Proton und Neutron.
\end{examplebox}

\begin{examplebox}
Betrachte eine horizontale Linie mit
\[
    Y=0.
\]
Dann gilt:
\[
    Q=I_3.
\]

Fuer ein Isospin-Triplett mit
\[
    I_3=+1,0,-1
\]
folgen die Ladungen:
\[
    Q=+1,0,-1.
\]

Das passt zum Pion-Triplett oder zum Sigma-Triplett.
\end{examplebox}

\section{Warum gibt es Oktett, Dekuplett und Nonett?}

\begin{conceptbox}
Die Multipletts entstehen aus der Kombination der leichten Quarkflavors
\[
    u,\quad d,\quad s.
\]

Mesonen bestehen aus:
\[
    q\overline{q}.
\]

Baryonen bestehen aus:
\[
    qqq.
\]

Die moeglichen Kombinationen lassen sich nach Symmetrien ordnen.
Daraus entstehen Oktette, Dekupletts und Singuletts.
\end{conceptbox}

\begin{notebox}
Eine vollstaendige mathematische Beschreibung benutzt die Flavor-Symmetrie
\[
    SU(3).
\]
Fuer diese Vorlesung ist meist wichtiger, die Diagramme lesen und die Quantenzahlen
der Teilchen bestimmen zu koennen.
\end{notebox}

\begin{examtipbox}
Du musst nicht nur die Namen der Multipletts kennen, sondern auch die Logik:
\begin{itemize}
    \item Baryonenoktett: Spin-$\frac{1}{2}$-Baryonen.
    \item Baryonendekuplett: Spin-$\frac{3}{2}$-Baryonen.
    \item Mesonennonett: Mesonen aus $q\overline{q}$-Zustaenden.
\end{itemize}
\end{examtipbox}

\section{Mini-Zusammenfassung}

\begin{notebox}
Die wichtigsten Punkte dieses Kapitels sind:
\begin{itemize}
    \item Hadronenmultipletts ordnen Hadronen nach Quantenzahlen.
    \item Die typischen Achsen sind $I_3$ und $Y$.
    \item Die elektrische Ladung folgt aus:
    \[
        Q=I_3+\frac{Y}{2}.
    \]
    \item Fuer strange-Hadronen gilt:
    \[
        Y=B+S.
    \]
    \item Baryonen bestehen im einfachen Quarkmodell aus drei Quarks.
    \item Mesonen bestehen aus einem Quark und einem Antiquark.
    \item Proton:
    \[
        p=uud.
    \]
    \item Neutron:
    \[
        n=udd.
    \]
    \item Pionen bilden ein Isospin-Triplett.
    \item Kaonen enthalten ein strange-Quark oder strange-Antiquark.
    \item Das Baryonenoktett enthaelt Spin-$\frac{1}{2}$-Baryonen.
    \item Das Baryonendekuplett enthaelt Spin-$\frac{3}{2}$-Baryonen.
    \item Das Omega-Minus hat Quarkinhalt:
    \[
        \Omega^- = sss.
    \]
    \item Mehr strange-Quarks bedeuten typischerweise groessere Masse.
\end{itemize}
\end{notebox}

\section{Pruefungscheck}

\begin{examtipbox}
Nach diesem Kapitel solltest du folgende Fragen beantworten koennen:
\begin{enumerate}
    \item Was ist ein Hadronenmultiplett?
    \item Welche Achsen verwendet man haeufig in Multiplettdiagrammen?
    \item Wie lautet die Gell-Mann--Nishijima-Formel?
    \item Was ist Hyperladung?
    \item Wie berechnet man $Y$ aus $B$ und $S$?
    \item Was ist der Quarkinhalt von Proton und Neutron?
    \item Welche Pionen gibt es und welchen Quarkinhalt haben sie?
    \item Welche Kaonen und Antikaonen gibt es?
    \item Warum hat $K^+$ positive Strangeness?
    \item Was ist das Mesonenoktett beziehungsweise Mesonennonett?
    \item Welche Teilchen gehoeren zum Baryonenoktett?
    \item Welche Teilchen gehoeren zum Baryonendekuplett?
    \item Was ist der Quarkinhalt von $\Omega^-$?
    \item Warum unterscheiden sich Proton und $\Delta^+$ trotz gleichem Quarkinhalt?
    \item Wie berechnet man die Ladung eines Hadrons aus seinem Quarkinhalt?
    \item Warum steigen Massen innerhalb eines Multipletts oft mit der Zahl der strange-Quarks?
    \item Was bedeutet Spin-$\frac{1}{2}$ im Baryonenoktett?
    \item Was bedeutet Spin-$\frac{3}{2}$ im Baryonendekuplett?
\end{enumerate}
\end{examtipbox}